\chapter{General Conclusion}

\section{Summary of Work Completed}

This project gave birth to \textbf{\caimp{}}, an open-source infrastructure
monitoring platform built around a philosophy unprecedented in the monitoring
world: answering the question ``\textit{Why is my server broken?}'' directly,
rather than displaying graphs for the operator to interpret alone.

In the space of one month, a team of three engineering students designed,
developed, tested and documented a complete platform of 19 microservices, over
52,000 lines of code in Python 3.12, Go 1.22 and TypeScript, a relational
database of 22 tables integrating TimescaleDB and pgvector, and a modern
answers-first user interface. The integration of a local large language model
(Llama 3.1 8B via Ollama), a change correlation engine, and Holt-Winters
statistical forecasting places \caimp{} at the intersection of software
engineering, data science and applied artificial intelligence.

Key delivered features include: real-time anomaly detection by three
complementary algorithms; local LLM natural-language explanation; automatic
correlation of system changes with incidents; 24-hour forecasting with confidence
intervals; intelligent alert deduplication; PDF reports accessible to
non-technical stakeholders; and a microservices architecture secured by mTLS,
JWT and RLS PostgreSQL.

\section{Academic and Professional Contributions}

This project mobilised a set of cross-disciplinary skills rarely all called upon
in a single academic module: distributed architecture, advanced databases
(TimescaleDB, pgvector), applied AI (LLM integration, prompt engineering, RAG
pipeline), data engineering (Holt-Winters, time-series processing, correlation
scoring), full-stack security (PKI, mTLS, JWT, RBAC, RLS), and DevOps
(containerisation, versioned migrations, platform self-monitoring).

\section{Personal Reflections}

\paragraph{Rim CHAHIDY}
This project allowed me to consolidate my React frontend skills and discover
the fascinating complexity of distributed system design. Building the Telemetry
Writer in Go was a stimulating technical challenge that introduced me to the
language's native concurrency mechanisms. I take from this experience the
conviction that user interface quality can transform a functional tool into
a genuinely useful solution.

\paragraph{Wiam SAMDI}
Working on the RAG pipeline and LLM prompts was the most intellectually enriching
part of this project. Learning to communicate with a language model to obtain
structured and reliable responses is an art requiring rigour and extensive testing.
This project also confirmed the importance of integration tests in distributed
systems: the hardest bugs never appear in unit tests.

\paragraph{Youssef TAÏK}
Designing end-to-end security — from the internal certificate authority through
to PostgreSQL RLS policies — was the most formative aspect of my contribution.
This project taught me that good security architecture is not a layer you add
at the end, but a backbone you build from day one. The vision of \caimp{} as
a virtual SRE for small teams feels like a genuinely promising direction.

\section{Closing Remarks}

\caimp{} was born from a simple conviction: developers who manage their own
servers should not need to become monitoring experts to know whether their
systems are functioning. Technology is now sufficiently mature — local LLMs,
vector databases, lightweight event-driven architectures — to automate most
of the investigation work that a monitoring alert traditionally requires.

We hope this work will contribute, modestly, to democratising access to
intelligent infrastructure monitoring, and that it will serve as a solid
foundation for teams wishing to extend and improve it.

\bigskip
\begin{flushright}
\textit{Settat, 2026}
\end{flushright}
